\section{Introdução}
Com o crescimento exponencial da conectividade e da complexidade das infraestruturas de rede, a segurança cibernética tornou-se uma preocupação crítica para organizações e indivíduos. Os ataques evoluíram em sofisticação, utilizando técnicas mapeadas pelo \mitreattack e explorando vulnerabilidades criticamente ativas no mercado, como as catalogadas pelo CISA KEV, exigindo ferramentas eficazes de teste de penetração e análise de vulnerabilidades que permitam avaliar proativamente a resiliência dos sistemas.

Neste contexto, este projeto apresenta o desenvolvimento da ferramenta \offscan, uma solução ofensiva para redes Wi-Fi escrita em Go.

O \offscan foi desenvolvido com foco em modularidade, desempenho e autonomia, reduzindo as dependências externas por meio do uso da linguagem Go, que possui bibliotecas robustas para redes de computadores. A ferramenta incorpora funcionalidades para testes ofensivos, incluindo descoberta de dispositivos e redes, captura e injeção de pacotes, e técnicas específicas para redes sem fio, como ataques a protocolos e criptografia.

Este relatório documenta o processo de desenvolvimento do \offscan, detalhando a evolução da arquitetura, as funcionalidades implementadas e os resultados obtidos em ambiente controlado. Além disso, discute as lições aprendidas, as limitações identificadas e as direções futuras para o projeto.